\section{Divisibility along the chain}\label{sec:chain-divisibility}

Fix the mesh point, $T$ and $u_0$ of \S\ref{sec:theta-chain}, an
auxiliary word $u$, and an inert prime $p$ as in \S\ref{sec:inert-divisibility}.
Write $K_j=K_{j,u}$ and $\mathcal K_j=K_j\cdots K_1$, with
$\mathcal K_0=1$. The bounds of Corollary~\ref{cor:inert-chain-jets}
add across lifts once the output lattice of each lift is the input
lattice of the next.

\subsection{From the kernel to integral forms}

Let $\mathcal V$ be the places $v\mid p$ of $E^+$ split in $E$.
Write $g_{A,vj},g_{B,vj}$ for \eqref{eq:inert-base-gain} and
\eqref{eq:inert-parity-gain} at lift $j$, and
$\alpha_v=\sum_{s\ {\rm even}}\alpha_s$ at $v$. This last sum is
independent of $j$ by \eqref{eq:chain-weights}. Use a discretely
valued coefficient field in $\mathbb C_p$ with valuation ring
$\mathcal O$ and uniformizer $\pi$, where $v_p(\pi)=1/e$.
For a lattice $H$ put $p^bH=\pi^{\lceil eb\rceil}H$. Then
$p^a(p^bH)\subseteq p^{a+b}H$.

At level $j$ write $H_j(L)$ for the coefficient lattice in the frame
of $L$, and $I_j$ for the forms with $F(L)\in H_j(L)$ at every $L$.
The model uses products of minors of the highest weight at even
positions and of its dual at odd positions, with the determinant
twists of \citep[Definition~11.3]{adriazola-grh}. Fix $1\le r\le k/2$.
At $v$ a Hecke edge of type $i\le r$ replaces $L_w$ by $\Lambda$ with
$pL_w\subseteq\Lambda\subseteq L_w$ and
$\dim_{\kappa(w)}L_w/\Lambda=i$, and takes the dual lattice at $\bar w$.
Let $T^{\rm raw,(j)}_{v,i}$ sum over these edges. Type zero is the identity.
Suppress $v$ in $\lambda_s=\lambda^{(0)}_{\tau_s}$ and $w_j(s)=w_j(\tau_s)$.
For $j\ge1$ put
\[
 n^{(j)}_{v,i}=\sum_{s\ {\rm even}}\sum_{a=1}^{i}\lambda_s(a)
                   +i\sum_{s\ {\rm odd}}w_j(s),\qquad
 T^{(j)}_{v,i}=p^{n^{(j)}_{v,i}}T^{\rm raw,(j)}_{v,i}.
\]

For a choice of types $\mathbf i=(i_v)_v$, let $L''$ be reached from
$L'$ by these edges. In the frames of $L,L''$, assign to each minor
of height $h$ and column set $I$ in \S\ref{sec:inert-divisibility} the exponent
\[
 \begin{cases}
 \min(h,i_v)-|I\cap[i_v]|,&s\ \text{even},\\
 |I\cap[i_v]|,&s\ \text{odd}.
 \end{cases}
\]
Here $[0]=\varnothing$. Write $E_v$ for the sum of these exponents at $v$.

\begin{hypothesis}[\openinput]\label{hyp:chain-lattice-transfer}
The lifts at $T,u_0$ have a common coefficient model with equivariant
lattices $H_j(L)$. For every $j,u,L,L',L''$ as above, use the prescribed
kernel of Hypothesis~\ref{hyp:inert-kernel-integrality} in the adapted
frames. If every frame coefficient $c$ of every product of jets in
Corollary~\ref{cor:inert-chain-jets} satisfies $v_p(c)+\sum_v E_v\ge b$,
then the map of \eqref{eq:chain-class-sum} satisfies
\begin{equation}\label{eq:chain-entry-inclusion}
 p^{\sum_v n^{(j)}_{v,i_v}} A_j(L'',L)H_{j-1}(L)
          \subseteq p^b H_j(L').
\end{equation}
\end{hypothesis}

At even positions Cauchy--Binet gives the transfer from bounds for
the individual row minors. Their relation to the frame coefficients
is examined in \S\ref{sec:integral-lifts}. At odd positions the transfer
uses the complementary minors in the dual
representation and the determinant factor
$(\det g/\det g')^{-w_j(s)}$ for the two changes of frame
\citep[Theorem~11.35, Lemma~11.37 and Hypothesis~11.38]{adriazola-grh}.

Set $d_j=\max_{[L]}v_p(|\operatorname{Stab}_{U_{j-1}}(L)|)$ over the
source classes in \eqref{eq:chain-class-sum}, and put
\[
 b_{vj}(0)=\max\{g_{A,vj},g_{B,vj}-\alpha_v\},\qquad
 b_{vj}(i)=g_{B,vj}\quad(1\le i\le r).
\]
\begin{lemma}\label{lem:chain-one-step}
Assume Hypotheses~\ref{hyp:inert-kernel-integrality}
and~\ref{hyp:chain-lattice-transfer}. For each choice of types,
\[
 \Big(\prod_v T^{(j)}_{v,i_v}\Big)K_j I_{j-1}
       \subseteq p^{\sum_v b_{vj}(i_v)-d_j} I_j .
\]
In particular $\mathcal K_m I_0\subseteq
p^{\sum_{j,v}b_{vj}(0)-\sum_j d_j}I_m$.
\end{lemma}

\begin{proof}
For $i_v>0$ take $J=[i_v]$. An even height-one top contributes one
to $E_v$ exactly when its index lies outside $J$. The other even
contributions are nonnegative, and all odd contributions are
$|I\cap J|$. Hence $E_v\ge N_J$. For $i_v=0$ we have $E_v=0$ and
$N_\varnothing=\alpha_v$. Corollary~\ref{cor:inert-chain-jets} now
gives $v_p(c)+\sum_v E_v\ge\sum_v b_{vj}(i_v)$ for this same product
of jets. Apply \eqref{eq:chain-entry-inclusion} to each term of the
edge sums and \eqref{eq:chain-class-sum}. Each term has one stabilizer
denominator. Taking every type zero and iterating gives the last inclusion.
\end{proof}

\subsection{Transport of a Hecke operator}

At level zero put
\[
 A_v=\sum_{s\ {\rm even}}\sum_{a=1}^r\lambda_s(a)
        -\sum_{s\ {\rm odd}}\sum_{a=1}^r\lambda_s(k+1-a),
 \qquad y_v=p^{A_v}T^{\rm raw,(0)}_{v,r}.
\]
The $y_v$ commute since their edges are at distinct places. In the
stated lattice model $y_v I_0\subseteq I_0$ by the edge inclusion
of \citep[Lemmas~11.4 and~11.8]{adriazola-grh}.
Write $q_v=p^{\nu_v}$, so $\#\kappa(w)=q_v^2$, and set
\[
 D_{vj}(X)=\sum_{i=0}^r q_v^{-i(N_j-i)}T^{\rm raw,(j)}_{v,i}X^i .
\]

\begin{hypothesis}[\openinput]\label{hyp:chain-hecke}
For the same word and frames, at every lift and orbit there are
$\epsilon_{vj}\in\mathcal O^\times$ and $\eta_{vj}$ with
$v_p(\eta_{vj})=\sum_{s\ {\rm even}}w_j(s)-\sum_{s\ {\rm odd}}w_j(s)$
such that, modulo $X^{r+1}$,
\begin{equation}\label{eq:chain-one-hecke}
 D_{vj}(X)K_j =
 \prod_{a=0}^{t_r}(1+\epsilon_{vj}\eta_{vj}q_v^{t_r-2a}X)
 \,K_jD_{v,j-1}(\epsilon_{vj}X).
\end{equation}
\end{hypothesis}

The split-lattice count gives \eqref{eq:chain-one-hecke} once the
frame and determinant factors of $A_j$ vary with the two lattices by
the prescribed length factors \citep[Proposition~10.13 and
Hypothesis~11.22]{adriazola-grh}. Hypothesis~\ref{hyp:inert-kernel-integrality}
controls the series for each pair separately.

\begin{lemma}\label{lem:chain-hecke-price}
Assume Hypothesis~\ref{hyp:chain-hecke}. There is an operator
$Z_{vj}=\sum_{i=0}^r c_i T^{(j)}_{v,i}$ with $c_i\in\mathcal O$ and
$v_p(c_0)\ge\alpha_v$, such that
\[
 Z_{vj}\mathcal K_j=\rho_{vj}\mathcal K_j y_v,\qquad
 v_p(\rho_{vj})=\delta_{vj}
      :=n^{(j)}_{v,r}-A_v+\nu_v r(N_j-N_0).
\]
\end{lemma}

\begin{proof}
Fix $v$ and suppress it. Write $P_j(X)$ for the scalar product in
\eqref{eq:chain-one-hecke}, set $Q_0=1$,
$Q_j(X)=P_j(X)Q_{j-1}(\epsilon_jX)$, and put $U_j=\prod_{\ell=1}^j\epsilon_\ell$.
Iteration gives
$D_j(X)\mathcal K_j=Q_j(X)\mathcal K_jD_0(U_jX)$ modulo $X^{r+1}$.
Since $Q_j(0)=1$, multiply by $Q_j^{-1}$ and take the coefficient of $X^r$.
After multiplication by $p^{n^{(j)}_r}q^{r(N_j-r)}$ the coefficients are
\[
 c_i=([X^{r-i}]Q_j^{-1})\,p^{n^{(j)}_r-n^{(j)}_i}
              q^{r(N_j-r)-i(N_j-i)},\qquad
 \rho_j=U_j^r p^{n^{(j)}_r-A}q^{r(N_j-N_0)} .
\]
The coefficient estimate of \citep[Theorem~11.33]{adriazola-grh}
gives $v_p(c_i)\ge0$ and $v_p(c_0)\ge\alpha_v$. Its hypotheses are
$2r\le k$, the weight bounds and increasing nodes of \S\ref{sec:theta-chain},
and the displayed valuation of $\eta_{vj}$ in Hypothesis~\ref{hyp:chain-hecke}.
The formula for $\delta_{vj}$ follows since $U_j$ is a unit.
\end{proof}

\subsection{The composite}

Let $B\subseteq\mathcal V$ and $y_C=\prod_{v\in C}y_v$, with
$y_\varnothing=1$. Put
\[
 h_{vj}=(g_{B,vj}-\delta_{vj}-g_{A,vj})_+,\qquad
 S_v=\sum_{j=1}^m h_{vj},\qquad
 D_B=\sum_{C\subseteq B} y_C^m
       \bigl(p^{\sum_{v\in B\setminus C}S_v} I_0\bigr).
\]
Thus an orbit in $C$ supplies a Hecke factor for each lift. An orbit
in $B\setminus C$ supplies the corresponding divisibility at the source.

\begin{theorem}\label{thm:chain-divisibility}
Assume Hypotheses~\ref{hyp:inert-kernel-integrality},
\ref{hyp:chain-lattice-transfer} and~\ref{hyp:chain-hecke}. At the
chosen $T,u_0$,
\[
 \mathcal K_m D_B\subseteq p^{G_B} I_m,\qquad
 G_B=\sum_{j=1}^m\left(\sum_{v\in\mathcal V}g_{A,vj}-d_j\right)
          +\sum_{v\in B}S_v .
\]
\end{theorem}

\begin{proof}
For $V\subseteq\mathcal V$, expand $\prod_{v\in V} Z_{vj}$.
At $v\in V$ a positive edge gives $g_{B,vj}$ by Lemma~\ref{lem:chain-one-step}.
The identity edge gives at least $g_{B,vj}-\alpha_v$, and its coefficient
supplies $\alpha_v$ by Lemma~\ref{lem:chain-hecke-price}. Thus
\[
 \Big(\prod_{v\in V}Z_{vj}\Big)K_j I_{j-1}
 \subseteq p^{\sum_{v\in V}g_{B,vj}+\sum_{v\notin V}g_{A,vj}-d_j}I_j .
\]
Fix $C\subseteq B$. Set $V_j=\{v\in C:h_{vj}>0\}$ and
$a_j=\sum_v g_{A,vj}-d_j+\sum_{v\in C}h_{vj}$.
We prove $\mathcal K_j y_C^j I_0\subseteq p^{\sum_{\ell\le j}a_\ell}I_j$
by induction from $j=0$. For $F\in I_0$ put
$W=y_C^{j-1}y_{C\setminus V_j}F$. The last factor preserves $I_0$.
Commuting $y_{V_j}$ through the composite gives
\[
 \mathcal K_j y_C^j F
   =\Big(\prod_{v\in V_j}\rho_{vj}^{-1}\Big)
     \Big(\prod_{v\in V_j}Z_{vj}\Big)K_j\mathcal K_{j-1}W .
\]
The induction hypothesis applies to $\mathcal K_{j-1}W$. The last
lift adds $\sum_{v\in V_j}(g_{B,vj}-\delta_{vj})
+\sum_{v\notin V_j}g_{A,vj}-d_j=a_j$. This proves the induction.
For the summand indexed by $C$ in $D_B$, the initial scalar adds
$\sum_{v\in B\setminus C}S_v$, giving $G_B$ in every summand.
\end{proof}

For the period map of \S\ref{sec:theta-chain}, fix a coefficient place
$\mathfrak p\mid p$ and local loss allowances $a,b\ge0$.

\begin{hypothesis}[\openinput]\label{hyp:chain-endpoints}
At the chosen shift the source lattice and final pairing satisfy
\[
 N_{u,\mathfrak p}\subseteq p^{-a}D_B,\qquad
 P_u(I_m)\subseteq p^{-b}R^\vee_{\mathfrak p}.
\]
\end{hypothesis}

Under Hypothesis~\ref{hyp:chain-endpoints}, Theorem~\ref{thm:chain-divisibility}
gives $B_u(N_{u,\mathfrak p})\subseteq p^{G_B-a-b}R^\vee_{\mathfrak p}$
by $B_u=P_u\mathcal K_m|_{N_u}$. Thus the local bound for the period
map depends on the two lattice losses at this same $u_0$.
